\documentclass[10pt,a4paper]{amsart}
\usepackage[margin=19mm]{geometry}
\usepackage{amsmath,amssymb,mathtools}
\usepackage[scr=rsfs,cal=euler]{mathalfa}
\usepackage{xfrac}
\usepackage[unicode,hidelinks,bookmarks=false]{hyperref}
\newcommand{\ie}{i.e.\ }
\newcommand{\cf}{cf.\ }
\newcommand{\resp}{respectively\ }
\newcommand{\Subsection}{\subsection}
\providecommand{\nobreakdash}{\nobreak\hbox{-}}
\setlength{\emergencystretch}{2em}
\multlinegap0pt
\usepackage{bm}
\usepackage{bbm}
\usepackage{dsfont}
\usepackage{xspace}
\usepackage{cancel}
\usepackage{mathtools}
\usepackage{amsthm}
\makeatletter
\@ifundefined{theorem}{\newtheorem{theorem}{Theorem}}{}
\@ifundefined{lemma}{\newtheorem{lemma}[theorem]{Lemma}}{}
\makeatother
\title[Finitary Cartesian closed varieties and semigroup actions]{Finitary Cartesian closed varieties and semigroup actions}
\author{Mark V Lawson}
\date{Source: arXiv:2602.19904v1 (CC BY 4.0)}
\begin{document}
\maketitle

\noindent\emph{Source and licence.} Finitary Cartesian closed varieties and semigroup actions. Version arXiv:2602.19904v1 (CC BY 4.0), distributed under the Creative Commons Attribution 4.0 International licence, as recorded in the arXiv licence metadata for this exact version and shown on the abstract page https://arxiv.org/abs/2602.19904v1. Full rights notice: licenses/arxiv-2602.19904-CC-BY-4.0.txt.\par\smallskip

\noindent\textit{Editorial scope.} Counted unit: the Theorem whose source label is \texttt{them:marytwo} in the article (source0.tex lines 862--864, source label \texttt{them:marytwo}), delivered together with the article's own complete original proof of it (source0.tex lines 865--985, one \texttt{proof} environment of 121 lines). No mathematics is written, repaired or re-derived: statement and proof are the article's own text line for line. Required context retained uncounted: them:mary (lines 516--518); lem:anne (lines 699--701). Every definition, display and label of those lines is retained unchanged. Omitted from this package and not redistributed: 1--515, 519--698, 702--861, 986--1370. Nothing inside a retained excerpt is omitted or altered: every retained body is delivered exactly as the article's lines are written, so no omission receipt of any kind accompanies this package. Citations: the retained excerpts carry no citation command at all, so no bibliography entry of the article is transcribed and no reference is redistributed. Reference adaptation: none was needed and none was made; every reference inside the delivered unit resolves inside it, so no unresolved reference or citation is produced. Printed slips kept literal: no printed word, symbol, hypothesis or number of the counted statement or of its proof was altered. Layout and class: the article's own class is \texttt{amsart} with packages including url, amsmath, cancel, graphicx, xy, eucal, amssymb, amsfonts, stackengine, xcolor; the wrapper is the lane's pinned \texttt{amsart} editorial wrapper at 10pt on A4, which restates the theorem-like environments the retained text needs and adds the article's own macro definitions that the retained text needs (none), with extra wrapper packages (bm, bbm, dsfont, xspace, cancel, mathtools). The delivered numbering is the unit's own. Assets: no retained span contains a figure environment, a graphics-inclusion command, a PDF-page inclusion, a table or a TikZ picture; no image file is copied into this package, the package declares no asset dependency and no image file is redistributed with it. Licence: version arXiv:2602.19904v1 of this article is distributed under the Creative Commons Attribution 4.0 International licence, as recorded in the arXiv licence metadata for this exact version and shown on the abstract page https://arxiv.org/abs/2602.19904v1. A full rights notice is delivered at licenses/arxiv-2602.19904-CC-BY-4.0.txt. Characters: every retained excerpt is pure ASCII, so the delivered engine, XeLaTeX with its default Latin Modern text font, reports no missing character.
\begin{theorem}\label{them:mary} Let $S$ be a left restriction monoid.  
The category $S$-{\bf FSup} is equivalent to the category $[\mathsf{Proj}(S)|\mathsf{Tot}(S)]$-{\bf Set}.
\end{theorem}

\begin{lemma}\label{lem:anne} Let $S$ be a Boolean left restriction monoid.
Then $[\mathsf{Proj}(S)|\mathsf{Tot}(S)]$ is a matched pair (in the above sense).
\end{lemma}

\begin{theorem}\label{them:marytwo} Let $S$ be a Boolean left restriction monoid.  
The category $S$-{\bf FSup} is equivalent to the category $[\mathsf{Proj}(S)|\mathsf{Tot}(S)]$-{\bf Set}.
\end{theorem}

\begin{proof} Put $M = \mathsf{Tot}(S)$, $B = \mathsf{Proj}(S)$. 
We first define a functor from $S$-{\bf FSup} to $[B|M]$-{\bf Set}.
In fact, $[B|M]$ is a matched pair by Lemma~\ref{lem:anne}.
Suppose $(S,X,p)$ is a supported action.
Put $X_{1} = p^{-1}(1)$.
Then $X_{1}$ is a $[B|M]$-set by Theorem~\ref{them:mary}.
It just remains to check that the additional axioms hold.
It is straightforward to show that (MPA7) and (MPA8) hold.
To show that axiom (MPA9) holds,
let $e \in B$ and $x,y \in X_{1}$.
Recall that $x \equiv_{e} y$ iff $e \cdot x = e \cdot y$.
Observe that
$e \cdot x \approx \bar{e} \cdot y$.
Thus $z = e \cdot x \vee \bar{e} \cdot x$ is well-defined.
But $z \equiv_{e} x$ and $z \equiv_{\bar{e}} y$.
Thus (MPA9) really holds.
Let $(S,X,p)$ and $(S,X',p')$ be supported actions
and
let $\theta \colon X \rightarrow X'$ be a homomorphism of supported actions.
The proof that $\theta | X_{1} \colon X_{1} \rightarrow X'_{1}$ defines a homomorphism of $[B|M]$-sets
is now routine.
It is now easy to see that we have defined a functor from the category of supported actions of $S$
to the category {\bf $[B|M]$-sets}.

We now define a functor from $[B|M]$-{\bf Set} to  $S$-{\bf FSup}.
This will follow Theorem~\ref{them:mary}.
Let $Y$ be a $[B|M]$-set.
We assume that $Y$ is non-empty.
As before, 
define
$$X = \bigcup_{e \in \mathsf{Proj}(S)} Y/\equiv_{e}$$
and
$p \colon X \rightarrow \mathsf{Proj}(S)$
by $p([x]_{e}) = e$.
Define 
$$s \bullet [x]_{e} = [\hat{s} \cdot x]_{(se)^{+}}.$$
This defines a supported action of $S$ on $X$, but we now need to show that the axioms
(E3)--(E7) hold for $X$.

(E3) holds: 
Put $z = [y]_{0}$ where $y \in Y$ is any element.
By (MPA7), it does not matter which element $y$ we choose.
From the definition of the action, this really is the minimum element of the poset $X$.
Again from the definition of the action, we have proved that this axiom holds.

(E4) holds: clear.

(E5) holds.
We need a preliminary result first:
if $x \equiv_{ef} y$ then there exists a $w \in X$ such that
$w \equiv_{e} x$ and $w \equiv_{f} y$.
Such a $w$ is unique upto $\equiv_{ef}$-equivalence.
By (MPA9)
there is a $w \in X$ such that $w \equiv_{e} x$ and $w \equiv_{\bar{e}} y$.
In a Boolean algebra, we have that $f = ef \vee \bar{e}f$.
By (MPA4),
we have that $w \equiv_{ef} x$ and $w \equiv_{\bar{e}f} y$. 
It follows that $w \equiv_{ef} y$ and $w \equiv_{\bar{e}f} y$.
By (MPA8), we have that that $w \equiv_{f} y$.
Now, suppose that there exists a $w' \in X$ such that
$w' \equiv_{e} x$ and $w' \equiv_{f} y$.
Then $w \equiv_{e} w'$.
But $ef \leq e$.
Thus by (MPA4),
we have that  $w \equiv_{ef} w'$, as required.

We can now prove that axiom (E5) holds.
Suppose that  $[x]_{e} \approx [y]_{f}$.
We prove that
$[w]_{e + f} = [x]_{e} \vee [y]_{f}$
where $w \in X$ is any element such that 
$w \equiv_{e} x$ and $w \equiv_{f} y$.
Observe that $x \equiv_{ef} y$ means precisely $[x]_{e} \approx [y]_{f}$.
Thus if $[x]_{e} \approx [y]_{f}$, then, using the element $w$ guaranteed above,
we have that $[x]_{e}, [y]_{f} \leq [w]_{e + f}$.
Suppose that $[x]_{e}, [y]_{f} \leq [w']_{i}$.
Then $e \leq i$ and $x \equiv_{e} w'$, and $f \leq i$ and $y \equiv_{f} w'$.
Thus $e + f \leq i$, and $w \equiv_{e} z$ and $w \equiv_{f} z$.
By (MP8), $w' \equiv_{e + f f} z$.
It follows that $[w]_{e + f} \leq [w']_{i}$.
We have therefore proved that $[w]_{e + f} = [x]_{e} \vee [y]_{f}$.

(E6) holds.
We show that if $[x]_{e} \approx [y]_{f}$ then 
$s \bullet ([x]_{e} \vee [y]_{f}) = s \bullet [x]_{e} \vee s \bullet [y]_{f}$.
Let $s = s^{+}\hat{s}$.
Then $s \bullet ([x]_{e} \vee [y]_{f}) = [\hat{s} \cdot z]_{(s(e + f))^{+}}$.
Also, $s \bullet [x]_{e} = [\hat{s} \cdot x]_{(se)^{+}}$ and $s \bullet [y]_{f} = [\hat{s} \cdot y]_{(sf)^{+}}$.
From $z \equiv_{e} x$ and $z \equiv_{f} y$ and using (MPA6),
we have that $\hat{s} \cdot z \equiv_{\hat{s} \ast e} \hat{s} \cdot x$ and $\hat{s} \cdot z \equiv_{\hat{s} \ast f} \hat{s} \cdot y$.
Observe that $(se)^{+} \leq (\hat{s}e)^{+}$ and $(sf)^{+} \leq (\hat{s}f)^{+}$.
Thus by (MPA4), we have that 
$\hat{s} \cdot z \equiv_{(se)^{+}} \hat{s} \cdot x$ and $\hat{s} \cdot z \equiv_{(sf)^{+}} \hat{s} \cdot y$.
The result now follows, from the following calculation
$$(s(e+f))^{+} = (se \vee sf)^{+} = (se)^{+} + (sf)^{+}.$$

(E7) holds.
Let $s \vee t = (s \vee t)^{+}\hat{u}$ where $\hat{u} \in M$.
Thus $s \vee t = s^{+}\hat{u} \vee t^{+}\hat{u}$.
But $s = s^{+}\hat{u}$ and $t = t^{+}\hat{u}$.
The result now follows.

Let $X$ and $X'$ be $[B|M]$-sets.
Let $\alpha \colon X \rightarrow X'$ be a homomorphism of $[B|M]$-sets.
Put $Y = \bigcup_{e \in \mathsf{Proj}(S)} X/\equiv_{e}$
and
$Y' = \bigcup_{e \in \mathsf{Proj}(S)} X'/\equiv_{e}$.
Define $\theta \colon Y \rightarrow Y'$ by $\theta ([x]_{e}) = [\alpha (x)]_{e}$.
Then this is a homomorphism of supported actions by Theorem~\ref{them:mary}.
However, we also have to check that $\theta$ maps the smallest element of $Y$ to the smallest element of $Y'$
and that binary joins in $Y$ are preserved by $\theta$.
The proofs of these are both straightforward from the definition of $\alpha$.
We have therefore defined a functor from $[B|M]$-{\bf Set} to  $S$-{\bf FSup}.

To prove that we have an equivalence of categories, it is enough to prove that 
every supported action of a Boolean left restriction monoid is isomorphic to an induced supported action.
By Theorem~\ref{them:mary} there is an isomorphism,
but we need to check that minimum elements are mapped
and binary joins are preserved.
Both of these are easy.
\end{proof}

\end{document}
